\documentclass[10pt,a4paper]{amsart}
\usepackage[margin=19mm]{geometry}
\usepackage{amsmath,amssymb,mathtools}
\usepackage[scr=rsfs,cal=euler]{mathalfa}
\usepackage{xfrac}
\usepackage[unicode,hidelinks,bookmarks=false]{hyperref}
\newcommand{\ie}{i.e.\ }
\newcommand{\cf}{cf.\ }
\newcommand{\resp}{respectively\ }
\newcommand{\Subsection}{\subsection}
\providecommand{\nobreakdash}{\nobreak\hbox{-}}
\setlength{\emergencystretch}{2em}
\multlinegap0pt
\usepackage{amssymb}
\usepackage{bbm}
\usepackage{enumerate}
\usepackage{cancel}
\usepackage{bm}
\usepackage{xcolor}
\newtheorem{assumption}{Assumption}
\newtheorem{proposition}{Proposition}
\newtheorem{lemma}{Lemma}
\newcommand{\C}{\mathbf{C}}
\newcommand{\E}{\mathbf{E}}
\newcommand{\Lin}{\mathcal{L}}
\newcommand{\N}{\mathbf{N}}
\newcommand{\R}{\mathbf{R}}
\newcommand{\Var}{\mathrm{Var}}
\newcommand{\X}{\mathcal{X}}
\title{The eigenvalues of i.i.d. matrices are hyperuniform}
\author{Giorgio Cipolloni, L\'aszl\'o Erd\H{o}s, Oleksii Kolupaiev}
\date{Source: arXiv:2602.17628v2 (CC BY 4.0)}
\begin{document}
\maketitle

\noindent\emph{Source and licence.} The eigenvalues of i.i.d. matrices are hyperuniform. Version arXiv:2602.17628v2 (CC BY 4.0), distributed by arXiv under the Creative Commons Attribution 4.0 International licence, http://creativecommons.org/licenses/by/4.0/, as recorded in the arXiv licence metadata for this exact version and shown on the abstract page https://arxiv.org/abs/2602.17628v2. Full licence notice: \texttt{licenses/arxiv-2602.17628-CC-BY-4.0.txt}.\par\smallskip

\noindent\textit{Editorial scope.} Counted unit: the article's lemma lem:E\_- (source0.tex lines 2330--2336). The article's complete original proof of it is delivered with it (source0.tex lines 2338--2378, one \texttt{proof} environment of 41 lines). No mathematics is written, repaired or re-derived: the statement and the proof are the article's own lines, in the article's own notation, and no retained line is substituted or re-flowed.  Retained uncounted as the source-local context the proof needs: lines 342--355; lines 366--372; lines 455--458; lines 468--471; lines 482--495; lines 506--512; lines 587--603.  Nothing was omitted from any retained span, so the package carries no omission record.  No retained line contains a citation, so the wrapper carries no bibliography and no bibliography entry.  No retained line contains a figure environment, an image-inclusion command or drawing code, so no image is delivered and the package declares no asset dependency.  Editorial formatting only, in addition: an AMS \texttt{amsart} preamble that restates the article's own declarations and macros used by the retained lines, namely the environments \texttt{assumption}, \texttt{proposition}, \texttt{lemma} on separate counters, together with the standard packages those lines need.  The retained environments print in this document as Assumption 1, Proposition 1, Lemma 1 in order of appearance, whereas the article prints the same items with the numbering of its own full document.  The complete native source of the article as distributed in the arXiv e-print for this exact version is bundled unchanged as \texttt{sources/194976/source0.tex}.
\begin{assumption}\label{ass:chi} \emph{(i)} The random variable $\chi$ satisfies $\E \chi=0$, $\E |\chi|^2=1$. In addition, in the complex case we also assume $\E\chi^2=0$. We further assume the existence of high moments, i.e. that there exist constants $C_p>0$, for any $p\in \N$, such that  
\begin{equation}
\E |\chi|^p\le C_p.
\label{eq:momass}
\end{equation}

\noindent\emph{(ii)} There exist a (small) $\mathfrak{a}>0$ and a (large) $\mathfrak{b}>0$ such that the probability density $\rho_\chi$ of $\chi$
%denoted by $\rho_\chi$,
satisfies
\begin{equation}
\rho_\chi \in L^{1+\mathfrak{a}}(\C),\quad \|\rho_\chi\|_{1+\mathfrak{a}}\le N^\mathfrak{b}.
\label{eq:reg_ass}
\end{equation}
\end{assumption}

\begin{assumption}\label{ass:Omega} The domain $\Omega_N\subset\C$ is open and simply connected, with $C^2$ boundary. There exists $\delta>0$ such that $\Omega_N\subset (1-\delta)\mathbf{D}$ for all $N\in\N$. There further exists an exponent $\alpha\in [0,1/2)$, such that $\Omega_N=N^{-\alpha}\widetilde{\Omega}_N$, where 
\begin{equation}
{\rm{diam}}(\widetilde{\Omega}_N):=\sup \{|z_1-z_2|\,:\, z_1,z_2\in\widetilde{\Omega}_N\} \sim 1,
\label{eq:diam}
\end{equation}
and the curvature of the boundary $\partial\widetilde{\Omega}_N$ is bounded from above uniformly in $N$.
\end{assumption}

\begin{equation}
f_{a,N}:=\phi_N*\omega_{a,N},
\label{eq:def_f}
\end{equation}

\begin{equation}
\omega_{a,N}^{(z_0)}(z):=\omega_{a,N}(z-z_0) = N^{2a}\omega(N^a(z-z_0)),\quad z,z_0\in\C,
\label{eq:omega_z0}
\end{equation}

\begin{proposition}\label{prop:E_omega} Let $X$ be a complex $N\times N$ i.i.d. matrix satisfying Assumption \ref{ass:chi}. Fix $\delta>0$, $a\in [1/2,1/2+\nu_0)$ with $\nu_0:=1/14$, and let $\omega\in C^2_c(\C)$ be a complex-valued function with integral 1. Uniformly in $z_0\in\mathbf{D}$ with $|z_0|\le 1-\delta$ it holds that
\begin{equation}
\E\Lin_N\big(\omega_{a,N}^{(z_0)}\big) = \frac{N}{\pi}\left(1+\mathcal{O}(N^{-c})\right)
\label{eq:E_omega}
\end{equation}
for some $c>0$, where $\omega_{a,N}^{(z_0)}$ is defined in \eqref{eq:omega_z0}.

Fix additionally $\alpha\in [0,1/2)$. Let $\Omega_N\subset \mathbf{D}$ be an $N$-dependent domain satisfying Assumption~\ref{ass:Omega} with exponent $\alpha$ and let $f_{a,N}$ be defined as in~\eqref{eq:def_f} with $\phi_N$ equal to the test function of $\Omega_N$. Then we have
\begin{equation}
\E\Lin_N\big(f_{a,N}\big) = \frac{N|\Omega_N|}{\pi}\left(1+\mathcal{O}(N^{-c})\right),
\label{eq:E_f}
\end{equation}
where $|\Omega_N|$ is the volume of $\Omega_N$.
\end{proposition}

\begin{proposition}\label{prop:Var_main} Assume the set-up and conditions of Proposition \ref{prop:E_omega}, but instead of $a\in [1/2,1/2+\nu_0)$ consider any fixed $a\ge 1/2$. Then it holds that
\begin{equation}
\Var\left[ \Lin_N(f_{a,N})\right]\lesssim N^{2(a-\alpha) - 2q_0(1/2-\alpha)+\xi} 
\label{eq:Var_main}
\end{equation}
for any fixed $\xi>0$, where $q_0:=1/20$.
\end{proposition}

\subsection{Ideas behind the proofs of Propositions \ref{prop:E_omega} and \ref{prop:Var_main}}\label{sec:results_list} For any $z\in\C$ denote
\begin{equation}
H^z:=W-Z,\quad \text{where}\quad 
W=\begin{pmatrix}
0&X\\X^*&0
\end{pmatrix}\in \C^{(2N)\times(2N)},\quad
Z=\begin{pmatrix}
0&z\\\overline{z}&0
\end{pmatrix}\in \C^{(2N)\times(2N)}.
\label{eq:def_hermitization}
\end{equation}
Here $Z$ is a $2\times 2$ block-constant matrix with blocks of size $N\times N$. The matrix $H^z$ is called the \emph{Hermitization} of $X$, while $z$ is known the \emph{Hermitization parameter}. We denote the resolvent of $H^z$ by
\begin{equation}
G^z(w):=(H^z-w)^{-1}=(W-Z-w)^{-1},\quad w\in \C\setminus\R,
\label{eq:def_resolvent}
\end{equation}
and call $w$ a \emph{spectral parameter}.

\begin{lemma}\label{lem:E_-}  For any $n\in\N$ and for $G=G^z(w)$ with $z\in\C$ and $w\in\C\setminus\R$, it holds that
\begin{align}
\langle (GE_-)^{2n-1}\rangle &= 0,\label{eq:E_-_identities1}\\
\langle (GE_-)^{2n}\rangle &= 2\sum_{s=0}^{n-1} (-1)^{s+1} \binom{2n-2-s}{n-1} \frac{\langle G^{s+1}\rangle}{(2w)^{2n-s-1}}. \label{eq:E_-_identities2}
\end{align}
Moreover, \eqref{eq:E_-_identities1}--\eqref{eq:E_-_identities2} remain valid after taking the deterministic approximations to all quantities appearing in these identities.
\end{lemma}

\begin{proof}[Proof of Lemma \ref{lem:E_-}] We start with the proof of~\eqref{eq:E_-_identities1}. Inverting the $2\times 2$-block matrix $W-Z-w$ defined in~\eqref{eq:def_hermitization} we get
\begin{equation}
G\!=\!\begin{pmatrix}
w \widetilde{H}_z(w)& X_z H_z(w)\\ H_z(w)X_z^*& w H_z(w)
\end{pmatrix}\!\!,\,\, \text{where}\,\, X_z\!:=X-z,\,
H_z(w)\!:=\! (X_z^*X_z-w^2)^{-1}\!\!,\, \widetilde{H}_z(w)\!:=\! (X_zX_z^*-w^2)^{-1}\!\!.
\label{eq:G_expl}
\end{equation}
Here the matrices $X_z^*X_z-w^2$ and $X_zX_z^*-w^2$ are invertible, since either $\Re w=0$, in which case $-w^2=(\Im w)^2>0$, or $\Re w\neq 0$, which implies $\Im (w^2)\neq 0$. Denote additionally $\widetilde{G}:=G^z(-w)$. From~\eqref{eq:G_expl} and the symmetry relations $\widetilde{H}_z(w)=\widetilde{H}_z(-w)$ and $H_z(w)=H_z(-w)$ we have that
\begin{equation}
GE_-=-E_-\widetilde{G}.
\label{eq:GE}
\end{equation}
Thus, the lhs. of~\eqref{eq:E_-_identities1} can be written as
\begin{equation}
\begin{split}
\langle (GE_-)^{2n-1}\rangle &= \langle (GE_-GE_-)^{n-1} GE_-\rangle = (-1)^{n-1} \langle (G\widetilde{G})^{n-1} GE_-\rangle\\
& = (-1)^{n-1} (2w)^{-n+1}\langle (G-\widetilde{G})^{n-1} GE_-\rangle=w\langle (\widetilde{H}_z(w))^{n}-(H_z(w))^{n}\rangle.
\end{split}
\label{eq:odd_k}
\end{equation}
To go from the first to the second line we employed the resolvent identity and in the last identity we used~\eqref{eq:G_expl}. Observing additionally that $\widetilde{H}_z(\eta)$ and $H_z(\eta)$ have identical spectra, we conclude that the rhs. of~\eqref{eq:odd_k} equals to zero. This finishes the proof of~\eqref{eq:E_-_identities1}.

To prove~\eqref{eq:E_-_identities2}, we compute similarly to~\eqref{eq:odd_k} that
\begin{equation}
\langle (GE_-)^{2n}\rangle = (-1)^n \langle (G\widetilde{G})^n\rangle =(-1)^n \langle G^n\widetilde{G}^n\rangle = -\frac{\partial_{w_1}^{n-1}\partial_{w_2}^{n-1} \langle G^z(w_1)G^z(-w_2)\rangle\big|_{w_1=w_2=w}}{((n-1)!)^2}.
\label{eq:even_k}
\end{equation}
Combining the resolvent identity with \eqref{eq:GE} we get
\begin{equation}
\langle G^z(w_1)G^z(-w_2)\rangle = \frac{\langle G^z(w_1)\rangle - \langle G^z(-w_2)\rangle}{w_1+w_2} = \frac{\langle G^z(w_1)\rangle}{w_1+w_2} + \frac{\langle G^z(w_2)\rangle}{w_1+w_2}.
\label{eq:even_k_res_id}
\end{equation}
Finally, we differentiate the first and the second term in the rhs. of~\eqref{eq:even_k_res_id} $n-1$ times in $w_2$ and $w_1$, respectively, and combine the result with \eqref{eq:even_k}:
\begin{equation*}
\langle (GE_-)^{2n}\rangle = \frac{(-1)^n2}{(n-1)!}\partial_y^{n-1} \frac{\langle G^z(y)\rangle}{(y+w)^n}\big|_{y=w}= \frac{2}{(n-1)!} \sum_{s=0}^{n-1}(-1)^{s+1} \frac{(2n-2-s)!}{(n-1-s)!} \cdot \frac{\langle G^{s+1}\rangle}{(2w)^{2n-1-s}}.
\end{equation*}
This finishes the proof of \eqref{eq:E_-_identities2}.

The proof of the analogues of~\eqref{eq:E_-_identities1} and~\eqref{eq:E_-_identities2} for the deterministic approximations follows by a standard meta-argument and thus is omitted.
\end{proof}

\end{document}
